
\chapter{Stress or Strain Sensor}
Stress analysis is the evaluation of the stress distribution within a body, focusing on determining the type, magnitude, and direction of the stresses involved. For complex geometries, theoretical stress analysis becomes less reliable because many assumptions must be made, also the mathematics become highly complicated. In such cases, experimental stress analysis is used \cite[see][75]{Sirohi:1991}. Direct measurement of mechanical stress is not possible, so strain is typically measured and subsequently converted into the corresponding stress values using established relationships \cite[see][16]{Murray:1992}. To understand the use of these relationships, it is important to first understand what strain is:  

When a force is applied to an object, it causes a deformation that depends on the object's shape, the material it is made of, and the way the force is applied. Figure \ref{fig:ForceOnBar}  illustrates one of the simplest possible systems: a force acting on a body in uniaxial tension. In this case, the resulting deformation is a straightforward axial elongation \cite[see][1]{Holister:1967}.

\Mynote{Add the Equations to visualize the Relationship}

\begin{figure}[H]
	\centering
	\includegraphics[width=0.5\linewidth]{../../../Images/Syde/HTA/StrainGauge/ForceOnBar}
	\caption[Simple System of Loading]{Simple System of Loading\cite[see][1]{Holister:1967}}
	\label{fig:ForceOnBar}
\end{figure}

\section{Methods to measure Strain}
In the market, a variety of different strain gauges can be found. In addition to mechanical methods, opto-mechanical methods and electrical strain gauges that are discribed here, there are also measurement techniques and methods such as the grid method, the moiré fringe technique, the interferometry, and the photoelasticity to be considered \cite[see][75]{Sirohi:1991}. 

\subsection{Mechanical Methods}
Before the 1930s, all measurements of deformation were conducted using mechanical instruments. Mechanical components such as wedges, screws, levers, gears, and combinations thereof were employed to amplify the measured signals. Given their typical length of approximately 250 mm, these gauges recorded the average deformation over relatively large sections, making them less suitable for detecting steep strain gradients. Measurement accuracy was influenced by factors such as friction, backlash, weight and inertia effects, and the flexibility of the mechanical parts and Dynamic measurements were not feasible with such devices. The figure below  \ref{fig:MechanicalStrainGauge}  illustrates a mechanical strain gauge, where a lever amplifies the strain and transfers the deformation to a dial gauge for readout\cite[see][76]{Sirohi:1991}.

\begin{figure}[H]
	\centering
	\includegraphics[width=0.5\linewidth]{../../../Images/Syde/HTA/StrainGauge/MechanicalStrainGauge}
	\caption[Mechanical Strain Gauge]{Mechanical Strain Gauge\cite[see][76]{Sirohi:1991}}
	\label{fig:MechanicalStrainGauge}
\end{figure}

\subsection{Opto-mechanical Methods}
In contrast, opto-mechanical methods for strain measurement use optical levers for magnification, allowing the gauges to be much shorter in length, typically around 50 mm. An example of this type of strain gauge is the Tuckerman gauge \ref{fig:TuckermanGauge} . It is an optical-mechanical extensometer that measures strain using a combination of a right-angle roof prism and a tungsten-carbide rocker, known as the Lozenge. When the specimen deforms, for example, by elongation, the edge of the rocker tilts outward, causing its mirror surface to rotate by a small angle $\phi$. A collimated beam of light, provided by an auto-collimator, is reflected by the prism-rocker assembly. The tilt of the rocker deflects the reflected beam by an angle $2\phi$ resulting in a linear shift of the image by $2f\phi$, with f as the focal lenght, at the focal plane of the auto-collimator. This shift, proportional to both the tilt angle and the focal length of the auto-collimator, is calibrated to correspond directly to the strain in the specimen. An important advantage of the Tuckerman gauge is that it does not require a fixed spatial relationship between the auto-collimator and the extensometer, which simplifies its practical use\cite[see][77]{Sirohi:1991}. 

\begin{figure}[H]
	\centering
	\includegraphics[width=0.5\linewidth]{../../../Images/Syde/HTA/StrainGauge/TuckermanGauge}
	\caption[Tuckerman Gauge]{Tuckerman Gauge\cite[see][77]{Sirohi:1991}}
	\label{fig:TuckermanGauge}
\end{figure}

\subsection{Electrical Methods}
On the other hand, electrical methods measure deformation by utilizing changes in specific electrical quantities. These methods are categorized into capacitive gauges, inductive gauges, piezoelectric gauges, and resistance gauges. Their main advantages are high sensitivity and the ability to respond to dynamic strains.

Inductive strain gauges detect deformation through variations in inductance, which result from changes in air gaps within the sensor\ref{fig:InductiveStrainGauge}. Capacitive methods, for example used in torque meters\ref{fig:CapacativeStrainGauge}, measure shifts in the relative positions of teeth under load, leading to changes in the effective area and thus in capacitance \cite[see][78]{Sirohi:1991}.

\begin{figure}[H]
	\centering
	\includegraphics[width=0.5\linewidth]{../../../Images/Syde/HTA/StrainGauge/InductiveStrainGauge}
	\caption[Inductive Strain Gauge]{Inductive Strain Gauge\cite[see][78]{Sirohi:1991}}
	\label{fig:InductiveStrainGauge}
\end{figure}

\begin{figure}[H]
	\centering
	\includegraphics[width=0.5\linewidth]{../../../Images/Syde/HTA/StrainGauge/CapacativeStrainGauge}
	\caption[Capacative Strain Gauge]{Capacative Strain Gauge \cite[see][78]{Sirohi:1991}}
	\label{fig:CapacativeStrainGauge}
\end{figure}

While capacitive and inductive gauges are generally larger in size and typically reserved for specialized applications, piezoelectric strain gauges are primarily used for studying dynamic deformations. In piezoelectric sensors, mechanical deformation generates a change in electrical voltage, enabling the accurate recording of rapidly varying strain signals \cite[see][78]{Sirohi:1991}.

\subsubsection{Unbonded Electrical resistance Strain Gauge}

Electrical resistance strain gauges detect strain by measuring changes in electrical resistance. There are two main types: unbonded and bonded strain gauges. In the unbonded version, two to twelve loops of wire are connected between both a stationary and a movable structures \ref{fig:UnbondedElectricalStrainGauge}. These gauges are primarily used as components in force and pressure transducers, as well as in accelerometers, rather than as direct strain measurement devices \cite[see][79]{Sirohi:1991}.

\begin{figure}[H]
	\centering
	\includegraphics[width=0.5\linewidth]{../../../Images/Syde/HTA/StrainGauge/UnbondedElectricalStrainGauge}
	\caption[Unbonded Electrical Strain Gauge]{Unbonded Electrical Strain Gauge \cite[see][79]{Sirohi:1991}}
	\label{fig:UnbondedElectricalStrainGauge}
\end{figure}

The bonded electrical resistance strain gauge, on the other hand, is cemented directly onto the surface of the body and is therefore subjected to the same strain as the material itself. The figure below shows different types of bonded strain gauges. One is mounted on a paper carrier, another on a Bakelite base, and a third on a plastic or rubberized, strippable base, as often used for foil gauges \ref{fig:BondedElectricalStrainGauge}. In the case of foil gauges, the carrier is typically stripped away before the gauge is cemented directly onto the test object.

\begin{figure}[H]
	\centering
	\includegraphics[width=0.5\linewidth]{../../../Images/Syde/HTA/StrainGauge/BondedElectricalStrainGauge}
	\caption[Types of Bonded Electrical Strain Gauges]{Types of Bonded Electrical Strain Gauges \cite[see][80]{Sirohi:1991}}
	\label{fig:BondedElectricalStrainGauge}
\end{figure}

This overview shows, that the market offers a wide range of mechanical, optical, and electrical strain gauges and extensometers, each designed for different purposes and applications, with varying characteristics and individual strengths. It is reasonable to assume that a single strain gauge cannot fulfill all desired specifications at once, especially when demanding conditions such as high-temperature environments are involved. However, for general use, the electrical resistance strain gauge meets most requirements \cite[see][17,18]{Murray:1992} and is well suited for integration into the HTA system.  

\subsubsection{Bonded Electrical resistance Strain Gauge}

The approach for electrical resistance strain gauges is based on a key discovery made in 1856 by the British physicist Lord Kelvin, who observed that certain electrical conductors exhibit a change in electrical resistance when subjected to mechanical strain. This proportional relationship between strain and resistance forms the fundamental operating principle of electrical resistance strain gauges \cite[see][18]{Murray:1992}. 
Wire-shaped conductors have the characteristic that their electrical resistance is directly proportional to their length and inversely proportional to their cross-sectional area. This relationship is described by the formula 

$$R=\rho \cdot \frac{l}{A}$$

where $R$ is the resistance, $l$ is the length of the conductor, $A$ is the cross-sectional area, and $\rho$ (rho) is a material-dependent constant known as the electrical resistivity. This coefficient reflects the inherent resistance of a material to the flow of electric current and depends on both the material itself and the units used for length and area. The formula confirms that increasing the length of a wire increases its resistance \cite[see][231]{Riecke:1896}. 

Eighty years later, two individuals independently developed the concept of the strain gauge. On the West Coast of the United States, in California, E.E. Simmons conducted material studies as part of his work at the California Institute of Technology (Caltech), focusing on the behavior of metals under impact loading by analyzing stress-strain diagrams. In order to measure the impact force, he used a dynamometer containing thin constantan resistance wires, representing the first documented use of the principle underlying strain gauges.  

At the same time, on the East Coast of the United States at the Massachusetts Institute of Technology (MIT), Professor Arthur S. Ruge was conducting research on the behavior of water tanks during earthquakes. Due to the thin walls of the tanks, the resulting strain in their vibrating table experiments was minimal. At that time, all conventional mechanical and optical measurement systems failed to detect such small deformations. However, attaching isolated thin wires from a potentiometer to the tank walls led to a change in the wires’ electrical resistance and a corresponding voltage drop, which yielded reliable measurements of the strain in the tank wall material.

\begin{figure}[H]
	\centering
	\includegraphics[width=0.7\linewidth]{../../../Images/Syde/HTA/StrainGauge/RugeStrainGauge1938}
	\caption[Ruge Strain Gauge \textbf{a} with and \textbf{b} without mounting bracket ]{Ruge Strain Gauge \textbf{a} with and \textbf{b} without mounting bracket \cite[see][3]{Keil:2017}}
	\label{fig:RugeStrainGauge1938}
\end{figure}

When producing and selling the strain gauge and for more comfortable handling, he glued the wires to a paper carrier, used plastic pieces to secure the connecting leads in place, and employed a brass tension clamp during the installation process, which was removed afterward \ref{fig:RugeStrainGauge1938} \cite[see][1-3]{Keil:2017}. The system of today strain gauges fundamentally remains the same, but they have been continuously optimized over the years in terms of size, accuracy, and installation process. 

\subsection{Loadcells} \Mynote{Picture of Circuit in Loadcell}
When Electrical Strain Gauges are bonded to a metal structure, the assembly is known as a load cell. Load cells are typically made from materials such as phosphor-bronze, aluminum, or stainless steel, and they are manufactured in well-defined shapes. When a mechanical quantity like force, weight, acceleration, or torque is applied, it creates a strain in the metal structure. This strain is sensed by the strain gauges and can result from axial loading, bending, or shear \cite[see][403]{Padmanabhan:2000}.

Load cells are typically categorized into three main structural designs: bending beam, shear beam, and column types \ref{fig:LoadcellTypes}. Variations of these basic designs can be seen in different form factors, depending on the application’s physical requirements and the desired measurement range. Typically, bending beam load cells find use in simpler applications and are well-suited for lower capacity measurements. In contrast, shear beam load cells excel in applications that demand higher durability and resistance to side loads, making them ideal for more demanding or fatigue-rated scenarios. Finally, column load cells are engineered to handle higher capacities, and they are most commonly deployed in heavy-duty industrial environments\cite[see][403]{honeywell:2013}.

\begin{figure}[H]
	\centering
	\includegraphics[width=0.7\linewidth]{../../../Images/Syde/HTA/StrainGauge/LoadcellTypes}
	\caption[Variety of Load Cells]{Variety of Load Cells \cite[see][403]{honeywell:2013}}
	\label{fig:LoadcellTypes}
\end{figure}

\section{Signal Conditioning}
Signal conditioning is a fundamental part of any measurement system, ensuring that the sensor’s output is properly prepared for further processing or analysis. The raw signals from sensors often require adjustments because they might be too weak, noisy, nonlinear, or in the wrong format altogether,such as needing conversion from analog to digital or vice versa. Common tasks in signal conditioning include amplification, noise removal, linearization, and format conversion, whether changing resistance to current or scaling voltage signals appropriately.

When measuring with a strain gauge, two main types of signal conditioning are typically required. First, the strain gauge’s resistance change must be converted into a voltage signal that can be processed. This involves using a Wheatstone bridge circuit, which translates the small resistance variations into a corresponding out-of-balance voltage. Second, the resulting voltage signal often needs to be amplified to an appropriate level—millivolt signals from the strain gauge must be increased to volts if they are to be digitized by an analog-to-digital converter. Together, these steps ensure the signal is both in the correct form and at a suitable level for further processing and accurate measurement \cite[see][78-79]{Bolton:2019}.



\subsection{Wheatstone Bridge}
The Bonded Electrical resistance Strain Gauge used, exhibits a change in resistance when subjected to mechanical strain. To enable a microcontroller like the ESP32-S3 to read this change, the resistance variation must be converted into a corresponding voltage signal. For this purpose, we employ a Wheatstone bridge \ref{WheatstoneBridge}, which translates the resistance change into a measurable voltage difference.

\begin{figure}[H]
	\centering
	\includegraphics[width=0.7\linewidth]{../../../Images/Syde/HTA/StrainGauge/WheatstoneBridge}
	\caption[Wheatstone Bridge]{Wheatstone Bridge \cite[see][93]{Bolton:2019}}
	\label{fig:WheatstoneBridge}
\end{figure}

To understand the principle of the Wheatstone bridge, Kirchhoff’s laws as well as Ohm’s law are essential. These laws help clarify how voltages, currents, and resistances interact in the circuit. By applying them, the relationship between the resistors in the Wheatstone configuration becomes clear.

When analyzing the Wheatstone bridge, the following facts emerge:
\begin{itemize}
	\item The total current is equal to the sum of $I_{1}$ and $I_{2}$: $I_{s} = I_{1} + I_{2}$
	\item The supply voltage $V_{s}$ is equal to the sum of the voltage drop across $R_{1}$ and across $R_{2}$ as well as to the sum of the voltage drop across $R_{3}$ and across $R_{4}$: $V_{s} = V_{31,2} = V_{3,4}$ 
	\item The voltage drop across one resistor in the Wheatstone bridge results from its resistance relative to the total resistance of its series branch. It can be calculated using the voltage divider rule: $V_{1}= V_{s}* R_{1}/(R_{1}+R_{2})$
	\item The potential difference between points B and D can be calculated by subtracting $V_{D}$ from $V_{B}$. This is equivalent to subtracting the voltage drop across $R_{4}$ from the voltage drop across $R_{2}$
	\item The resulting formula for the potential therefore is: $V_{B} - V_{D} = V_{s}* R_{2}/(R_{1}+R_{2}) - V_{s}* R_{4}/(R_{3}+R_{4}) $
	\item which concludes to: $ V_{s}* (R_{2}/(R_{1}+R_{2}) - R_{4}/(R_{3}+R_{4}))$

\end{itemize}

\Mynote{Find a book that provides the Laws!!}

The primary goal of the Wheatstone bridge is to measure the potential difference between points B and D. This measurement allows us to draw conclusions about the relationship between the resistors, particularly when one of them is unknown or subject to change as in the usage of strain gauges.

The concept of a “balanced” Wheatstone bridge refers to the condition in which the potential difference between B and D is zero. In this state, the voltage across the bridge diagonal vanishes, indicating that the ratio of resistors in both branches is equal.

With the Wheatstone bridge being balanced a one more statements can be made: 
\begin{itemize}
	\item With $ V_{s} = 0$ follows $R_{2}/(R_{1}+R_{2}) - R_{4}/(R_{3}+R_{4}) = 0 $ akkording to Zero Property of Multiplication
	\item Hence: $R_{2}*(R_{3}+R_{4}) = R_{4}*(R_{1}+R_{2})$
	\item distributive law: $R_{2}*R_{3}+R_{2}*R_{4} = R_{4}*R_{1}+R_{4}*R_{2}$
	\item The result is an equation that describes the relationship among the four resistors in a Wheatstone bridge: $R_{1}/R_{2}=R_{3}/R_{4}$
\end{itemize} 



\subsubsection{Bridge configurations}
A Wheatstone bridge can be configured in three ways: quarter-bridge, half-bridge, and full-bridge. In the quarter-bridge \ref{QuarterBridge} configuration, three fixed resistors and one active strain-sensitive resistor form the bridge. The half-bridge \ref{HalfBridge} uses two active strain-sensitive resistors and two fixed resistors, enhancing measurement sensitivity and reducing temperature-related errors. The full-bridge configuration \ref{FullBridge}, where all four resistors are active strain-sensitive elements, provides the highest sensitivity and best temperature compensation. In both the half- and full-bridge configurations, the active resistors are aligned so that one experiences compression while the other experiences tension, allowing the system to measure both positive and negative pressures \cite[see][73]{Tandeske:1991}.

\begin{figure}[H]
	\centering
	\includegraphics[width=0.5\linewidth]{../../../Images/Syde/HTA/StrainGauge/QuarterBridge}
	\caption[Quarter Wheatstone Bridge]{Quarter Wheatstone Bridge \cite[see][72]{Tandeske:1991}}
	\label{fig:QuarterBridge}
\end{figure}

\begin{figure}[H]
	\centering
	\includegraphics[width=0.5\linewidth]{../../../Images/Syde/HTA/StrainGauge/HalfBridge}
	\caption[Half Wheatstone Bridge]{Half Wheatstone Bridge \cite[see][72]{Tandeske:1991}}
	\label{fig:HalfBridge}
\end{figure}

\begin{figure}[H]
	\centering
	\includegraphics[width=0.5\linewidth]{../../../Images/Syde/HTA/StrainGauge/FullBridge}
	\caption[Full Wheatstone Bridge]{Full Wheatstone Bridge \cite[see][72]{Tandeske:1991}}
	\label{fig:FullBridge}
\end{figure}

\subsection{Amplification}
An amplifier can be described as a system with an input and an output, whose purpose is to increase the magnitude of a signal \ref{Amplifier}. Its voltage gain is defined as the ratio of the output voltage to the input voltage, each measured relative to the earth. Additionally, the input impedance of an amplifier is calculated by dividing the input voltage by the input current, while the output impedance is the output voltage divided by the output current \cite[see][79]{Bolton:2019}.

\begin{figure}[H]
	\centering
	\includegraphics[width=0.7\linewidth]{../../../Images/Syde/HTA/StrainGauge/Amplifier}
	\caption[Working System of an Amplifier]{Working System of an Amplifier\cite[see][79]{Bolton:2019}}
	\label{fig:Amplifier}
\end{figure}

\subsubsection{Programmable Gain Amplifier (PGA)}

While traditional operational amplifiers are widely used for many general-purpose tasks, there are also specialty amplifiers based on operational amplifiers designed for more specific applications. These include programmable gain amplifiers (PGAs) \cite[see][121]{Jung:2005}. 

An operational amplifier is a differential input, single-ended output amplifier used with external feedback elements to perform various operations. On top of this basic amplifier concept, a programmable gain amplifier (PGA) adds an important feature: the ability to control its gain in a predictable and adjustable way. Since ADCs typically have full-scale input voltage ranges between 2 V and 10 V, it is important that the input signals are adjusted close to this range to achieve the ADC’s full precision. Sensors, however, generate output voltages over a wide range, which means that a fixed gain could either fail to amplify small signals sufficiently or cause large signals to saturate the ADC or amplifier.

The PGA is usually placed between the sensor and the ADC to ensure that the ADC receives signals within its optimal range. In some cases, additional signal conditioning may be required before or after the PGA, but generally, it is preferable to place the PGA first to minimize errors introduced by further signal conditioning \cite[see][151-152]{Jung:2005}.

\begin{figure}[H]
	\centering
	\includegraphics[width=0.7\linewidth]{../../../Images/Syde/HTA/StrainGauge/PGALogic}
	\caption[PGAs in data acquisition systems]{PGAs in data acquisition systems \cite[see][152]{Jung:2005}}
	\label{fig:FullBridge}
\end{figure}

\subsection{Analog to Digital conversion}
Analog circuits can already handle a wide range of signal processing tasks effectively. Many functions, such as audio amplification, filtering, and radio transmission, are straightforward enough that analog processing provides sufficient performance. However, when more advanced functionality is required, digital signal processing (DSP) becomes indispensable, offering critical extensions of capabilities that go beyond the reach of analog circuits. The advantages of digital signal processing include perfect storage of digitized signals, an effectively unlimited signal-to-noise ratio, the ability to perform complex calculations, and the flexibility to adapt algorithms to changing circumstances. To leverage these benefits, it is essential to convert analog signals into digital form at an early stage of the processing chain, and to do so with high precision.

In analog systems, signals exist as continuous variations with respect to a physical reference. In contrast, digital systems represent signals as discrete numerical values within the limited word width of the digital format. The analog-to-digital converter (ADC) continuously strives to provide an optimal digital representation of the incoming analog signal at any given moment, despite the inherent rounding errors that must be accepted during the conversion process. Digital signals differ fundamentally from analog signals because they are both sampled and quantized. Sampling means that digital signals only have defined values at discrete time intervals set by the sampling frequency. Quantization, on the other hand, introduces finite resolution, mapping each sample to one of a limited set of digital values, which introduces small errors in the representation \ref{fig:ADCFunctions} \cite[see][2]{Pelgrom:2010}.

\begin{figure}[H]
	\centering
	\includegraphics[width=0.7\linewidth]{../../../Images/Syde/HTA/StrainGauge/ADCFunctions}
	\caption[Functions of the analog-to-digital converter: sampling, quantizing and linking to a reference]{Functions of the analog-to-digital converter: sampling, quantizing and linking to a reference \cite[see][2]{Pelgrom:2010}}
	\label{fig:ADCFunctions}
\end{figure}

\section{\textcolor{red}{Filter}}
\section{Software to read Strain}
To enable the HX711 to perform its intended functions within the measurement system, a software domain is needed. Here, the Arduino IDE serves as the primary platform for writing and managing the code that configures and controls the HX711’s operation. This environment provides framework for programming the ESP32-S3 microcontroller, which acts as the central unit for data acquisition and processing.

\subsection{Libraries}
To enable communication between the ESP32-S3 and the HX711, the Arduino HX711 library by Rob Tillaart is used. Next to the HX711 a Library called "HX711\_MP" exists providing a different calibration method.

\subsubsection{Library \FILE{HX711}} 
The HX711 library by Rob Tillaart builds on the earlier work of Bogde’s library, expanding its functionality. This library introduces additional functions that provide more detailed information about the measurement process, helping to extract the most accurate readings possible from the HX711 ADC. A key feature of this library is its use of floating-point numbers (floats). This approach leverages the 23-bit mantissa of the IEEE754 float format, which aligns well with the 24-bit resolution of the HX711 ADC. Additionally, using floats has the advantage of reducing memory usage on platforms such as the Arduino UNO. It’s important to note that, while the HX711’s ADC technically has 24-bit resolution, real-world measurements are typically limited to 16 to 20 meaningful bits due to inherent noise and other factors. As a result, single measurements can reliably provide around four to five significant digits. To improve accuracy and reduce noise, it’s recommended to average multiple readings.

The library has seen significant updates in its development. Starting with version 0.3.0, a custom \PYTHON{shiftIn()} function was introduced to ensure proper timing of the clock pulses when reading the ADC. This addresses timing issues observed on some processors that could lead to incorrect readings or missed sign bits. Although this approach adds an extra 50–55 microseconds to the reading process, it ensures more reliable and accurate data across different hardware platforms. In version 0.4.0, the \PYTHON{begin()} function was updated to include a new parameter: fastProcessor. This option allows the clock pulses to be slowed down for processors that can switch I/O pins faster than the HX711’s specification of $ 0.4 \mu s$. This ensures that timing remains within the chip’s requirements, preventing errors. While the HX711 chip can operate at either 10 or 80 samples per second (SPS) depending on the logic level at the RATE pin, this library currently does not offer built-in support to switch between these modes. In most breakout board designs, the RATE pin is permanently connected to ground, limiting operation to 10SPS. If needed, users can manually control the RATE pin using an additional I/O pin and the \PYTHON{digitalWrite()} function \cite[see][151-152]{Tillaart:2025}.

\subsubsection{Library \FILE{HX711\_MP}}
The HX711\_MP library is an extension of the original HX711 library (version 0.3.5). Although it builds on the same foundational code, it introduces several important changes and is not directly compatible with the original library’s interface. Currently, HX711\_MP is still considered experimental and requires further testing. However, since it is based on proven and well-tested code, no significant issues are expected by the creator. 

The main improvement of the HX711\_MP library lies in its use of multi-point calibration. Unlike the original library, which assumes a simple linear relationship between the raw ADC readings and the actual weight, the multi-point calibration method in HX711\_MP allows for up to ten calibration points. This makes it possible to correct for non-linear behavior in the sensor readings. Interpolation between calibration points remains linear, ensuring smooth transitions within the calibration range. It’s important to note that the HX711\_MP library does not extrapolate beyond the calibration data points. This means that measurements falling outside the calibrated range are not adjusted or interpolated, which is a crucial consideration when setting up the calibration data \cite[see][151-152]{Tillaart:2024}.

\subsubsection{Installation}
For this project, the focus will be on the HX711 library. To integrate the library into your project, it can be installed directly through the Library Manager in the Arduino IDE. To do this, open the Library Manager from the “Sketch” or “Tools” menu and search for “HX711” \ref{HX711LibraryInstall} by Rob Tillaart. Once located, a click on “Install” adds the library to the Arduino environment.

\begin{figure}[H]
	\centering
	\includegraphics[width=0.5\linewidth]{../../../Images/Syde/HTA/HX711LibraryInstall}
	\caption[HX711 Library Installation]{HX711 Library Installation}
	\label{fig:HX711LibraryInstall}
\end{figure}

\subsubsection{Interface}
The HX711 library offers a variety of functions that serve as the interface to the HX711 ADC. These functions cover essential tasks such as initializing the communication, setting the gain for the differential input, reading measurement data, and managing timing and data retrieval. By providing this interface, the library simplifies the integration of weight measurement functionality into different projects.

The main workflow is divided into three key steps: First, the connection to the HX711 module is established using begin(dataPin, clockPin). Next, calibration is carried out to ensure accurate measurements. Finally, the core logic of the application is implemented, which typically includes data readings and calculations. How to calibrate the HX711 will be discussed in more detail later in the section titled Calibration.

\paragraph{\#include "HX711.h":}
The statement \textbf{\#include <HX711.h>} at the beginning of the code integrates the HX711 library into the project, making its functions and classes available for use.

\paragraph{Constructor:}
The HX711 library includes the constructor \PYTHON{HX711()} and the destructor \PYTHON{~HX711()} for initializing and cleaning up the HX711 object, as well as functions for managing its operation. The \PYTHON{void begin(uint8\_t dataPin, uint8\_t clockPin, bool fastProcessor = false) } function establishes the connection with the specified data and clock pins, with an optional fastProcessor parameter to ensure proper timing on faster systems. Additionally, the \PYTHON{reset()} function sets the internal state of the chip to its default conditions, performing a power cycle in the process.

\paragraph{isReady:}
Several functions are provided to ensure that new measurement data is available before reading from the sensor. The \PYTHON{is\_ready()} function offers a simple check to see if the load cell is ready to deliver a fresh measurement. If the code needs to pause until new data is available, the \PYTHON{wait\_ready()} function can be used, which repeatedly checks the sensor at defined intervals (defaulting to zero milliseconds). For cases a limit to the number of attempts to wait for data is wanted, the \PYTHON{wait\_ready\_retry()} function comes into play, specifying both the maximum number of retries and the delay between them. Alternatively, the \PYTHON{wait\_ready\_timeout()} function waits until new data is ready or until a maximum timeout has passed, which by default is set to 1000 milliseconds.

\paragraph{Read:}
The \PYTHON{Read- functions} process the measurement data. The \PYTHON{float read()} function performs a raw read of the data. To calculate the average of multiple measurements, float \PYTHON{read\_average(uint8\_t times = 10)} can be used, where the times parameter defines the number of reads, with a minimum of 1. For more robust data, \PYTHON{float read\_median(uint8\_t times = 7)} returns the median value of multiple reads, where times can range from 3 to 15, preferably odd numbers for better statistical behavior. The \PYTHON{float read\_medavg(uint8\_t times = 7)} function calculates the average of the "middle half" of several readings, using the same range and preference for odd numbers. To smooth the data over time, the \PYTHON{float read\_runavg(uint8\_t times = 7, float alpha = 0.5)} function provides a running average across times measurements, with the alpha parameter controlling the weighting factor (between 0 and 1) for the new data point, requiring times to be at least 1. Finally, the \PYTHON{uint32\_t last\_read()} function retrieves the timestamp in milliseconds of the most recent measurement.

\paragraph{Gain + channel:}
For selecting the gain and channel functions are offered by the HX711 Library, though it’s important to use these functions carefully as their reliability is not fully guaranteed. Users are recommended to consult the datasheet before configuring these settings.
The available gain constants, defined in the header file, are \PYTHON{HX711\_CHANNEL\_A\_GAIN\_128} (default, gain of 128 for Channel A), \PYTHON{HX711\_CHANNEL\_A\_GAIN\_64} (gain of 64 for Channel A), and \PYTHON{HX711\_CHANNEL\_B\_GAIN\_32} (fixed gain of 32 for Channel B). The \PYTHON{set\_gain(uint8\_t gain = 128, bool forced = false)} function allows setting these values. If an invalid gain value is provided, no changes are made. With forced set to false, the library will only update the gain and channel if it detects a difference; when set to true, the function will always attempt to set the new gain and channel, including performing a dummy read to ensure the next measurement reflects the updated setting. The \PYTHON{get\_gain()} function returns the currently configured gain. It’s worth noting that changing the gain or channel can take up to 400 milliseconds, as specified in the datasheet. Due to the potential for the HX711 to enter an unknown state after an MCU reset or unexpected reboot, it is strongly recommended to explicitly call \PYTHON{set\_gain()} in the program’s \PYTHON{setup()} function to ensure the device is properly initialized. In practical use, achieving reliable channel and gain selection may be more difficult than the datasheet suggests, so careful testing is advised.

\paragraph{Read mode:}
The HX711 library allows setting and retrieving the operational mode for reading data using \PYTHON{get\_value()} and indirectly via \PYTHON{get\_units()}. Several read modes are defined as constants in the header file, including \PYTHON{HX711\_RAW\_MODE}, \PYTHON{HX711\_AVERAGE\_MODE}, \PYTHON{HX711\_MEDIAN\_MODE}, \PYTHON{HX711\_MEDAVG\_MODE}, and \PYTHON{HX711\_RUNAVG\_MODE}. In both median-based modes (\PYTHON{HX711\_MEDIAN\_MODE} and \PYTHON{HX711\_MEDAVG\_MODE})only 3 to 15 samples are permitted to help manage memory use.

A range of dedicated functions is available to set these modes. The \PYTHON{set\_raw\_mode()} function configures the library to perform a single read, while \PYTHON{set\_average\_mode()} calculates the average of multiple readings. \PYTHON{set\_median\_mode()} determines the median value from a series of samples, and \PYTHON{set\_medavg\_mode()} calculates the average of the central half of the sample data. For a running average, \PYTHON{set\_runavg\_mode()} is available, with the default smoothing factor set to 0.5. To check the current mode, the \PYTHON{get\_mode()} function returns the active read mode, which is set to \PYTHON{HX711\_AVERAGE\_MODE} by default.

\paragraph{Get values:}
To retrieve measurement data that has been adjusted for offset and scaled to meaningful units \PYTHON{Get values} are integrated in the Library. In \PYTHON{HX711\_RAW\_MODE}, the times parameter is ignored, and a single call to \PYTHON{read()} is performed instead. The  \PYTHON{float get\_value(uint8\_t times = 1)} function reads a value corrected for offset, while \PYTHON{float get\_units(uint8\_t times = 1)} returns the measured value converted to the appropriate physical units based on the scale factor. The \PYTHON{bool set\_scale(float scale = 1.0)} function sets the scale factor, which typically is a positive number greater than 50, depending on the specific load cell being used. Note that in some specialized cases, the scale might be negative. If scale is zero, the function returns false to indicate an error. The \PYTHON{float get\_scale()} function retrieves the currently set scale factor. Similarly, \PYTHON{void set\_offset(int32\_t offset = 0)} sets the offset, while \PYTHON{int32\_t get\_offset()} returns the current offset value.

\paragraph{Power Management:}
The HX711 library includes two functions for managing the chip’s power state. The \PYTHON{power\_down()} function disables the HX711, blocking for 64 microseconds. The \PYTHON{power\_up()} function reactivates the HX711 and should reset it to default settings, although this reset behavior is not always guaranteed.

\paragraph{Pricing:}
To support pricing applications, such as weighing goods or educational exercises, the library also includes a set of functions for handling price calculations. These include \PYTHON{float get\_price(uint8\_t times = 1)}, \PYTHON{void set\_unit\_price(float price = 1.0)} and \PYTHON{float get\_unit\_price() }.

\paragraph{Millivolts:}
The library provides experimental support for using the HX711 as a millivolt meter. The idea is to apply zero millivolts to the input and call \PYTHON{tare(times)} to establish the zero point, then apply 20mV and call \PYTHON{calibrate\_scale(20000)} to map the reading to microvolts. If the scale is linear, the \PYTHON{get\_units(uint8\_t times = 1)} function can then report microvolts, allowing the HX711 to be repurposed for measuring any other linear physical quantities.

\paragraph{Multiple HX711 Modules:}
For systems with multiple HX711 modules, the simplest approach is to use separate DOUT and CLK lines for each. Alternatively, a multiplexer can manage the connections, but this requires additional wiring and code. Another option is to share the CLK line among modules, which has trade-offs: while it can work if all HX711s use the same gain and channel settings, it can lead to mismatches after power cycles or resets. Sharing data lines (DOUT) is not possible, as it risks causing short circuits.

\subsection{Calibration}
Calibration of the HX711 with a load cell involves a series of steps to ensure accurate and consistent measurements. The first step is to clear the scale of any weight and call the \PYTHON{tare()} function to determine and set the zero-point offset. This step establishes a baseline measurement that represents zero weight. Once the tare is complete, a known weight is placed on the scale, and the \PYTHON{calibrate\_scale(float weight)} function is called. This weight is typically given in grams, but any unit can be used. During this step, the scale factor that converts the raw ADC readings to meaningful weight units is calculated. It is important to store the determined offset and scale factor, so they can be reused in future measurements without recalibration. The units used in \PYTHON{calibrate\_scale()} will also be the units returned by \PYTHON{get\_units()} in later measurements.

Several functions support this calibration process. The \PYTHON{tare(uint8\_t times = 10)} function, for example, averages multiple measurements to establish an accurate zero-point offset. The offset itself can be retrieved with \PYTHON{get\_tare()}, which returns the offset multiplied by the scale factor—this differs from the raw offset that can be obtained with \PYTHON{get\_offset()}. The \PYTHON{tare\_set()} function helps confirm that a tare has been successfully established by checking whether the offset is non-zero, which is the case for all tested load cells. The \PYTHON{calibrate\_scale()} function also takes an average of several measurements to improve the accuracy of the calibration. Starting with version 0.6.0, the weight parameter in \PYTHON{calibrate\_scale()} has been defined as a float, making it easier to calibrate with different units (such as kilograms) and to support a much broader range of calibration weights. For example, very heavy load cells, like those rated for 500kg , can be calibrated with a weight defined in grams (e.g., 100,000g) for improved precision. It’s worth noting that \PYTHON{calibrate\_scale()} uses its own averaging method and does not consider the current read mode of the library. It can also accept negative values, such as measuring the lifting force of a balloon, which may apply a negative load to the sensor. 

The second part of the calibration process involves understanding the tare and how it relates to real-world usage. Even with no load, a load cell typically produces a non-zero raw output, so \PYTHON{get\_tare()} is used to measure and establish this zero reference point. While it’s common to tare the scale with no load at all, it is also possible to define a zero point with a fixed load in place,like a cup, dish, or other container. This allows the system to automatically subtract the weight of the container from subsequent measurements. However, it’s crucial to ensure that the total load (fixed plus variable) stays within the load cell’s maximum capacity. For instance, if a load cell can handle 1000g and the fixed load (like a cup) weighs 300g, it’s best not to calibrate with weights that exceed 700g. Ideally, calibration should use weights that are about 80–90\% of the load cell’s maximum capacity to ensure accurate performance in the expected measurement range. Additionally, it’s recommended to calibrate with a weight that is close to the actual measurement range of interest. For example, if the application involves weighing portions of 250g, using a calibration weight between 200g and 300g can provide more accurate results.

Temperature effects also play a role, so calibration should ideally be performed at the same temperature in which the actual measurements will take place. Internally, the weight calculation is based on the formula: $Weight = get\_scale() * raw + get\_tare()$.
By interpolating this formula with the measured raw data, it’s possible to create an ideal conversion graph and compare it to actual data to assess the accuracy of the load cell setup.

\section{Example Applications of Strain measurement}
The following presents a series of example applications that demonstrate how strain can be accurately measured in practice. 

\subsection{Hardware used in Applications}
For these applications, several hardware components are needed to ensure accurate data collection and analysis like a BFH350-3AA Strain Gauge, an Soldered\_100384 load cell sensor 10kg and a HX711 Load Cell Amplifier.

\subsubsection{BFH350-3AA Strain Gauge}
For this setup the BFH350-3AA strain gauge model \ref{StrainGaugeBFH3503AA} will be used. This strain gauge has dimensions of 7.2 × 4.1 mm \ref{StrainGaugeUsed} and operates with a resistance of 350 ohms ±1 ohm, while the base material consists of a phenolic-epoxy compound. \Mynote{cite amazon}

\begin{figure}[H]
	\centering
	\includegraphics[width=0.5\linewidth]{../../../Images/Syde/HTA/StrainGauge/StrainGaugeBFH3503AA}
	\caption[BFH350 3AA ]{BFH350 3AA }%cite Amazon
	\label{fig:StrainGaugeBFH3503AA} 
\end{figure}

\begin{figure}[H]
	\centering
	\includegraphics[width=0.5\linewidth]{../../../Images/Syde/HTA/StrainGauge/StrainGaugeUsed}
	\caption[Strain Gauge Used]{Used Strain Gauge}
	\label{fig:StrainGaugeUsed}
\end{figure}

\subsubsection{Soldered\_100384 load cell sensor 10kg}
This 12x13x75mm load cell measures the weight of an applied load up to 10kg. It operates within a voltage range of 5V to 10V, providing an output of $2 \pm 0.1 mV/V$. Internally, the load cell integrates two fixed and two variable resistors arranged in a Wheatstone bridge configuration and produces a small output voltage, which requires amplification. To read these small output signals, the load cell must be connected to an HX711 board\cite[see][1]{soldered:2023}.

\begin{figure}[H]
	\centering
	\includegraphics[width=0.5\linewidth]{../../../Images/Syde/HTA/StrainGauge/Loadcell}
	\caption[Soldered\_100384 load cell sensor 10kg]{Soldered\_100384 load cell sensor 10kg\cite[see][1]{soldered:2023}}
	\label{fig:LoadCell} 
\end{figure}

\subsubsection{Hx711 Load Cell Amplifier}
The HX711 (\ref{HX711front} and \ref{HX711back}) is a high-resolution 24-bit ADC specifically tailored for applications that rely on bridge sensors, such as those used in weighing systems and industrial automation. It incorporates a programmable gain amplifier (PGA) characterized by low noise and capable of handling differential signals sourced via an input multiplexer, which allows users to toggle between Channel A and Channel B for measurements.

\begin{figure}[H]
	\centering
	\includegraphics[width=0.5\linewidth]{../../../Images/Syde/HTA/Setup/HX711front}
	\caption[HX711 front]{HX711 front}
	\label{fig:HX711front}
\end{figure}

\begin{figure}[H]
	\centering
	\includegraphics[width=0.5\linewidth]{../../../Images/Syde/HTA/Setup/HX711back}
	\caption[HX711 back]{HX711 back}
	\label{fig:HX711back}
\end{figure}

When powered with a 5V analog supply through the AVDD pin, Channel A can be configured with gain factors of 128 or 64, suitable for capturing signals with full-scale differential voltage ranges of $\pm 20 mV or \pm 40mV$, respectively. Channel B, on the other hand, operates with a fixed gain of 32. If there’s no external voltage regulator present, the HX711 can make use of its integrated voltage regulation circuitry, simplifying the setup. For clocking, the HX711 supports its own internal oscillator as well as options for external crystal or clock sources, offering flexibility in timing configurations. An embedded power-on-reset system makes the digital communication streamlined, removing the need for extensive setup procedures. The device’s operation is managed through dedicated pins and supports features that enhance performance and integration flexibility. These include a quiet and adjustable PGA (with gains of 32, 64, and 128), dual differential input channels, and an integrated power regulator that can deliver stable voltage to both the load cell and the ADC itself. The internal oscillator can also be synchronized with an external crystal when tighter timing precision is required. For data throughput, users have the choice of output data rates—either 80 samples per second or 10 samples per second, depending on the level of detail and speed required. Moreover, the ADC is designed to reject interference at the standard power-line frequencies of 50Hz and 60Hz. The HX711 consumes less than 1.5mA during normal operation and less than $1 \mu A$ in power-down mode. It has a operating voltage range from 2.6V to 5.5V and a temperature tolerance from $-40 ^{\circ}$ to $-80 ^{\circ}$. Housed in a 16-pin SOP-16 package, the HX711 is often used in applications requiring precise and reliable load cell measurements.

\paragraph{Block Diagram HX711}
\begin{figure}[H]
	\centering
	\includegraphics[width=0.8\linewidth]{../../../Images/Syde/HTA/Setup/HX711BlockDiagramm}
	\caption[HX711 Block Diagramm]{HX711 Block Diagramm \Mynote{Datasheet HX711}}
	\label{fig:HX711 Block Diagramm}
\end{figure}

The block diagram in Figure \ref{HX711BlockDiagramm} represents a typical application of the HX711 24-bit ADC in a weighing scale system. At its core, the HX711 interfaces directly with a bridge sensor, such as a load cell, which produces differential output signals proportional to the weight or force being measured. The load cell connects to the differential analog input pins of the HX711. These are split into two channels: Channel A (INA+, INA-) and Channel B (INB+, INB-). Channel A has the programmable gain amplifier (PGA) with selectable gains of 128 or 64 when powered by a 5V AVDD supply. Channel B has a fixed gain of 32. The differential input signals are routed through an input multiplexer to the low-noise PGA, which boosts the small analog signals to levels suitable for conversion. The amplified signals are then digitized by the integrated 24-bit sigma-delta ADC. An analog supply regulator is built into the HX711. This regulator supplies stable power to both the analog front end of the ADC and the load cell itself, simplifying the overall power management. It uses external components (transistor S8550, resistors R1 and R2) to create a precise analog voltage for AVDD, derived from the VSUP supply input. The regulator’s output is monitored and stabilized via the VFB and BASE pins. To manage system timing, the HX711 features an internal oscillator that requires no external components. If tighter timing control is needed, external clock signals can be applied to the XI and X0 pins, with the internal oscillator disabling itself automatically in this case.

The digital interface section of the HX711 includes two key pins: DOUT for serial data output and PD\_SCK for both serial clock input and power-down control. Data retrieval from the ADC is accomplished by pulsing the PD\_SCK pin (25–27 pulses per conversion), which shifts out the 24-bit data word through the DOUT pin. The data is presented in two’s complement format, and the output data rate can be set to 10 or 80 samples per second via the RATE pin, depending on application needs. The power supply connections are shown at the top of the diagram: VSUP (2.7V-5.5V) provides the main supply voltage, shared with the microcontroller unit (MCU) that receives the converted digital data. The analog ground (AGND) and digital ground (GND) provide return paths for the analog and digital domains, respectively, to help minimize noise.
\Mynote{HX711 Datasheet Ressource}


\subsection{Loadcell connected to a HX711}
This measurement system combines a load cell with an HX711 Amplifier to achieve precise weight measurements. The load cell itself contains a fully integrated Wheatstone bridge, which converts the mechanical force applied to the load cell into a small differential voltage signal. This sensitive output signal cannot be read directly by a microcontroller, so the HX711 is used to amplify and digitize it \ref{LoadCellCircuit}.

\begin{figure}[H]
	\centering
	\includegraphics[width=0.5\linewidth]{../../../Images/Syde/HTA/Setup/LoadcellCircuit}
	\caption[Load Cell Circuit]{Load Cell Circuit}%cite Amazon
	\label{fig:LoadCellCircuit} 
\end{figure}



\subsubsection{Circuit}
The load cell used in this setup incorporates a complete Wheatstone bridge, designed to convert the applied force into a differential voltage signal. The color coding of the load cell cables is typically identified and noted directly on the HX711 breakout board. This ensures that each wire is correctly connected to the appropriate input channel of the HX711. The wire connections are outlined in the following table\ref{LoadCellConnections}.

\begin{table}[H]
\begin{center}
\caption[Connections with the Terminals of Load Cell, HX711 and ESP32]{Connections with the Terminals of Load Cell, HX711 and ESP32 \cite[see][145]{Roobert:2025}}
\label{LoadCellConnections}
\begin{tabular}{|c|c|c|c|}
	\hline
	\textbf{LOAD CELL WIRE COLOR}& \textbf{H/11 (indicates)} & \textbf{HX711(refers to)} & \textbf{Microcontroller} \\
	\hline
	RED (E+)& E+ & GND & GND \\
	\hline
	BLACK (E-)& E- & DT & GPIO \\
	\hline
	WHITE (A-)& A- & SCK & GPIO \\
	\hline
	GREEN (A+)& A+ & VCC & 5V \\
	\hline
\end{tabular}
\end{center}
\end{table}

\subsubsection{\FILE{HX711\_kitchen\_scale} Example Code}
An example sketch named \FILE{HX711\_kitchen\_scale} can be found within the Arduino IDE’s example codes \ref{HXKitchenScaleCode}. This example demonstrates both the calibration process and regular measurement readings, providing a practical template for integrating the HX711 library into projects. It showcases how to initialize the sensor, perform a tare and calibration, and then acquire accurate weight readings in a typical application.

\begin{figure}[H]
	\centering
	\includegraphics[width=0.5\linewidth]{../../../Images/Syde/HTA/HXKitchenScaleCode}
	\caption[\FILE{HX711\_kitchen\_scale} Example Code]{FILE{HX711\_kitchen\_scale} Example Code}
	\label{fig:HXKitchenScaleCode}
\end{figure}

The \FILE{HX711\_kitchen\_scale} example sketch begins by initializing communication with the HX711 using \PYTHON{begin()}, specifying the data and clock pins connected to the sensor. It then calls \PYTHON{tare()} to determine the zero offset when the scale is empty, establishing a reliable baseline. Next, the user is prompted to place a known weight on the scale, and \PYTHON{calibrate\_scale()} is called to calculate the scale factor based on this known reference. This scale factor is used to convert the raw sensor output into meaningful weight measurements. In the main loop, the sketch continuously calls \PYTHON{get\_units()} to retrieve the measured weight in the desired units and displays the results on the serial monitor. This setup demonstrates the typical workflow for preparing and using a load cell with the HX711 to obtain accurate and reliable weight measurements.

\definecolor{codegreen}{rgb}{0,0.6,0}
\definecolor{codegray}{rgb}{0.5,0.5,0.5}
\definecolor{codepurple}{rgb}{0.58,0,0.82}
\definecolor{backcolour}{rgb}{0.95,0.95,0.92}

\lstinputlisting[
caption= FILE{HX711\_kitchen\_scale} Example Code,
label={lst:HX711Example},
language=C++,
backgroundcolor=\color{backcolour},   
commentstyle=\color{codegreen},
keywordstyle=\color{magenta},
numberstyle=\tiny\color{codegray},
stringstyle=\color{codepurple},
basicstyle=\ttfamily\footnotesize,
breakatwhitespace=false,         
breaklines=true,                 
keepspaces=true,                 
numbers=left,       
numbersep=5pt,                  
showspaces=false,                
showstringspaces=false,
showtabs=false,                  
tabsize=2]{../../../Code/Syde/HTADemonstrator/Test/HXKitchenScale/HXKitchenScale.ino}


\subsection{Strain Gauge Beam}
To visualize the impact of load on a beam and to measure the resulting strain, a strain gauge is installed on the surface of a steel beam. To convert these small changes in resistance into measurable signals, a half-bridge Wheatstone circuit is soldered to the strain gauge. This half-bridge arrangement is directly connected to an HX711 module, which amplifies and digitizes the signals. The digitized data is then processed by an ESP32-S3 microcontroller, which communicates with the HX711 to collect and interpret the strain measurements.

\subsubsection{Installation of Bonded Electrical resistance Strain Gauge}
For the installation of two strain gauges on a steel beam, the guideline "Strain Gauge Installation on Steel" \cite[see][]{Brown:2016} by SEL was used as a reference . However, slightly different materials and equipment were employed. As illustrated in the following figure, 2011 Araldite – a two-component adhesive by Huntsman – Magic Tape by Scotch, and Neutral RTV Silicone by Visbella were utilized during the bonding process. Subsequently, the strain gauges were covered with Butyl Waterproof Tape by ADHES for additional protection.
The strain gauges used for the initial prototype were of model BF350-3AA by YOUMILE, which come with thin wires pre-soldered to the contact pads. 


\Mynote{cite Website SEL after Email Response}

\begin{enumerate}
	\item The intended location for the strain gauge should be clearly marked on the steel surface. An area measuring approximately 2,5 × 5 cm is recommended to ensure sufficient space for installation and handling.
	\item An angle grinder equipped with a flapper wheel is recommended to remove any mill scale from the steel surface if necessary. Care must be taken to avoid excessive grinding, as this could reduce the cross-sectional area of the steel and compromise its structural integrity.
	\item A grinder and a die grinder with a stone attachment are used to smooth the area designated for the strain gauge. The surface should be uniformly smooth but does not require a mirror-like finish \ref{fig:BodyPrepare}. As in the previous step, care must be taken to avoid removing excessive material, which could compromise the cross-sectional area of the steel. Alternatively, a 100-grit sanding disk mounted on a pneumatic tool can also be used for this purpose.
	vvv
	\begin{figure}[H]
		\centering
		\includegraphics[width=0.5\linewidth]{../../../Images/Syde/HTA/StrainGauge/BodyPrepare}
		\caption[Surface Prepare]{Surface Prepare}
		\label{fig:BodyPrepare}
	\end{figure}
	
	\item The steel surface must be thoroughly cleaned. For this purpose, acetone is used in combination with non-woven sponges. The sponge is moistened with acetone—direct application of acetone to the surface is avoided, as it evaporates rapidly. The prepared sponge is then used to wipe the surface in a single, consistent direction. A new sponge is used each time, until it remains visibly clean (white) after wiping. Care must be taken not to touch unpolished areas with the sponge, in order to prevent contamination or damage. As acetone is highly flammable, appropriate safety precautions must be observed. This includes wearing safety goggles and ensuring adequate ventilation to avoid inhalation of acetone vapors.
	\item The next step is to mark the exact location for the strain gauge on the surface using a pencil or marker \ref{fig:MarkSteel}. One straight line perpendicular to the beam should be visible, so that the two arrows of the strain gauge can later be aligned to it \ref{fig:Positioning}. To prevent any contamination, the surface is cleaned with acetone again afterwards.
	
	\begin{figure}[H]
		\centering
		\includegraphics[width=0.5\linewidth]{../../../Images/Syde/HTA/StrainGauge/MarkSteel}
		\caption[Mark Steel]{Mark Steel}
		\label{fig:MarkSteel}
	\end{figure}
	
	\item Before attaching the strain gauge, it is essential to ensure that the surface remains clean. If there has been a delay between the preparation steps and the installation, an additional cleaning with acetone is recommended, as dust and other contaminants can accumulate quickly. In cases where the delay exceeds one to two days, a fine sanding disc should be used to lightly refresh the surface and ensure it remains smooth and free of contamination. 
	
	\begin{itemize}
		\item The following figure serves to identify the top and bottom sides of the strain gauge \ref{fig:StrainGaugeTop}. In this case, the soldering points clearly indicate that the second strain gauge is correctly oriented. Additionally, the shiny appearance of the wires in the first strain gauge suggests the upper side of the Strain Gauge. It is crucial to avoid touching the bottom surface, as this is the bonding side. Any contact may lead to contamination and compromise the quality of the adhesion.
		
		\begin{figure}[H]
			\centering
			\includegraphics[width=0.5\linewidth]{../../../Images/Syde/HTA/StrainGauge/StrainGaugeTop}
			\caption[Strain Gauge Top and Bottom]{Strain Gauge 1 Bottom an 2 Top}
			\label{fig:StrainGaugeTop}
		\end{figure}
		
		\item When removing the strain gauge from its protective plastic backing, it is attached to a piece of transparent Magic tape on the top side. A transparent, low-adhesive tape is needed, which is why Magic Tape is used in this case. Ideally, strain gauge installation tape should be chosen. The tape allows safe handling and precise positioning without direct contact. The gauge should be properly centered on the tape. To facilitate handling and ensure sufficient flexibility during the subsequent installation steps, the tape should be at least 5 cm in length. If there is a delay before installation, the taped gauge can be temporarily placed on a clean glass surface to preserve its condition. By utilizing the tape, the gauge can be positioned and repositioned precisely at the previously marked location\ref{fig:Positioning}. When removing the tape from any surface, it must always be lifted at a steep angle. This minimizes the risk of damaging the strain gauge.
		
		\begin{figure}[H]
			\centering
			\includegraphics[width=0.5\linewidth]{../../../Images/Syde/HTA/StrainGauge/Positioning}
			\caption[Strain Gauge Positioning]{Strain Gauge Positioning}
			\label{fig:Positioning}
		\end{figure}
		
		\item In the next step, the tape is partially peeled back, starting from the end where the soldering pads of the strain gauge are located, until the entire bonding surface of the gauge is exposed. The tape is not removed completely to ensure that the position of the strain gauge remains fixed and does not shift during the bonding process.		
		\item A very small amount of adhesive is now applied between the strain gauge and the steel surface \ref{fig:GlueProcess}. Using a smooth and continuous sliding motion with the thumb, the strain gauge is carefully repositioned into its final location. Once in place, firm and steady pressure is applied with the thumb for approximately one minute, to press the gauge into the adhesive layer and to eliminate any trapped air bubbles, ensuring proper adhesion\ref{fig:Thumb}.
		
		\begin{figure}[H]
			\centering
			\includegraphics[width=0.5\linewidth]{../../../Images/Syde/HTA/StrainGauge/GlueProcess}
			\caption[Strain Gauge Glue Process]{Strain Gauge Glue Process}
			\label{fig:GlueProcess}
		\end{figure}
		
		\begin{figure}[H]
			\centering
			\includegraphics[width=0.5\linewidth]{../../../Images/Syde/HTA/StrainGauge/Thumb}
			\caption[Strain Gauge Thumb Pressure]{Strain Gauge Thumb Pressure}
			\label{fig:Thumb}
		\end{figure}
		
		\item After the curing time of the adhesive, the tape can be slowly and carefully removed to avoid damaging the strain gauge or attached wires \ref{fig:GluedStrainGauge}. Unbounded areas may now be identified by the presence of air bubbles or slight differences in color between bonded and unbounded regions. If such areas appear near the edges of the gauge, additional adhesive can be applied. In these cases as well, firm pressure must be applied to ensure proper contact and to allow the adhesive to fully bond.
		
		\begin{figure}[H]
			\centering
			\includegraphics[width=0.5\linewidth]{../../../Images/Syde/HTA/StrainGauge/GluedStrainGauge}
			\caption[Glued Strain Gauge]{Glued Strain Gauge}
			\label{fig:GluedStrainGauge}
		\end{figure}
		
		\item The strain gauge wires should be securely fixed to the steel surface to prevent them from touching each other or being accidentally pulled off or damaged. For this purpose, painter’s tape or a similarly lightly adhesive tape is recommended.  
		
	\end{itemize}
	\item A waterproofing and mechanical protection layer is applied over the strain gauge to shield the adhesive bond from degradation caused by atmospheric moisture. This protective coating ensures the long-term durability of the installation and safeguards the gauge against environmental and mechanical influences \ref{fig:Silikon}. For the protective layer, silicone is used.
	
	\begin{figure}[H]
		\centering
		\includegraphics[width=0.5\linewidth]{../../../Images/Syde/HTA/StrainGauge/Silikon}
		\caption[Strain Gauge Silikon as a protective Layer]{Silikon as a protective Layer}
		\label{fig:Silikon}
	\end{figure}
	
	\item In this case, a half Wheatstone bridge is required to properly measure the axial strain on the beam. Therefore, the procedure is repeated on the exact opposite side\ref{fig:BeamStrainGauge}. The concept of the Wheatstone bridge is explained in the following chapter.
	
	\begin{figure}[H]
		\centering
		\includegraphics[width=0.5\linewidth]{../../../Images/Syde/HTA/StrainGauge/BeamStrainGauge}
		\caption[Beam with two Strain Gauges]{Beam with two Strain Gauges}
		\label{fig:BeamStrainGauge}
	\end{figure}
	
	The strain gauges are now wired according to the following wiring diagram \ref{fig:StrainGaugeWiring}. It is important to ensure that the wires do not touch the metal of the beam in order to keep the system insulated. Therefore, a protective layer of silicone is applied around the soldered connections \ref{fig:SiliconeLayer}.
	
	\begin{figure}[H]
		\centering
		\includegraphics[width=0.5\linewidth]{../../../Images/Syde/HTA/StrainGauge/StrainGaugeWiring}
		\caption[Strain Gauge Wiring]{Strain Gauge Wiring}
		\label{fig:StrainGaugeWiring}
	\end{figure}
	
	\Mynote{Redo the Picture - Ressource Youtube}
	
	\begin{figure}[H]
		\centering
		\includegraphics[width=0.5\linewidth]{../../../Images/Syde/HTA/StrainGauge/SiliconeLayer}
		\caption[Silicone Layer around Strain Gauge Wiring]{Silicone Layer around Strain Gauge Wiring}
		\label{fig:SiliconeLayer}
	\end{figure}
	
	\item The electrical resistance of the strain gauge is measured to ensure that it is not unintentionally electrically connected to the steel. This quality control step is essential, as in the worst-case scenario, a faulty installation would require the strain gauge to be removed and replaced. Using a multimeter set to measure resistance, the reading between each pair of wires of the strain gauge wires should correspond to the nominal resistance of the gauge, which is 350 ohms in this case \ref{fig:WireResistance}, \ref{fig:MultimeterWire}. Additionally, the resistance between the steel and each individual wire should display an overload or infinite resistance\ref{fig:SteelResistance}, \ref{fig:MultimeterSteel}, confirming that there is no unintended electrical connection to the structure.
	
	
	
	\begin{figure}[H]
		\centering
		\includegraphics[width=0.5\linewidth]{../../../Images/Syde/HTA/StrainGauge/WireResistance}
		\caption[Resistance Measurement between P+ and one the of strain Gauge wires ]{Resistance Measuring between P+ and one the of strain Gauge wires}
		\label{fig:WireResistance}
	\end{figure}
	
	\begin{figure}[H]
		\centering
		\includegraphics[width=0.5\linewidth]{../../../Images/Syde/HTA/StrainGauge/MultimeterWire}
		\caption[Measurement reading between P+ and one the of strain Gauge wires]{Measuring reading between P+ and one the of strain Gauge wires}
		\label{fig:MultimeterWire}
	\end{figure}
	
	\begin{figure}[H]
		\centering
		\includegraphics[width=0.5\linewidth]{../../../Images/Syde/HTA/StrainGauge/SteelResistance}
		\caption[Resistance Measurement between any one of the of strain Gauge wires and the beam]{Resistance Measuring between any one of the of strain Gauge wires and the beam}
		\label{fig:SteelResistance}
	\end{figure}
	
	\begin{figure}[H]
		\centering
		\includegraphics[width=0.5\linewidth]{../../../Images/Syde/HTA/StrainGauge/MultimeterSteel}
		\caption[Measurement reading between any one of the of strain Gauge wires and the beam]{Measurement reading between any one of the of strain Gauge wires and the beam}
		\label{fig:MultimeterSteel}
	\end{figure}
	
	After confirming that the strain gauges have been installed correctly and that no unintended contact with the metal occurs, two resistors with a value close to or identical to the nominal resistance of the strain gauges can be connected to the strain gauges within a Wheatstone bridge circuit, and then connected to the HX711 amplifier \ref{fig:WheatstoneAmp}. 
	
	\begin{figure}[H]
		\centering
		\includegraphics[width=0.5\linewidth]{../../../Images/Syde/HTA/StrainGauge/Wheatstone}
		\caption[Wheatstone Bridge connected to Amplifier]{Wheatstone Bridge connected to Amplifier}
		\label{fig:WheatstoneAmp}
	\end{figure}
	
\end{enumerate}

\subsubsection{Circuit}
The circuit for this application is essentially the same as for a Load Cell described earlier and follows the connections described in Table \ref{LoadCellConnections}. In this setup, however, the Wheatstone bridge is formed during the installation process, making it essential to carefully follow the wiring diagram for the Wheatstone bridge. Ensuring that the A+ and A-, as well as the E+ and E- connections, are correctly aligned with the HX711 module is critical for obtaining accurate measurements. It is important to note that if the positive and negative connections are reversed, the sign of the measurement readings will also be inverted. 

\subsubsection{Code}
This code sets up a simple web server on an ESP32-S3 microcontroller to remotely control and read weight data from the strain gauges. It allows users to tare, calibrate, and measure the load cell’s readings directly through a web interface, making it easy to manage and visualize the measurement process.

\definecolor{codegreen}{rgb}{0,0.6,0}
\definecolor{codegray}{rgb}{0.5,0.5,0.5}
\definecolor{codepurple}{rgb}{0.58,0,0.82}
\definecolor{backcolour}{rgb}{0.95,0.95,0.92}

\lstinputlisting[firstline=1,lastline=70,
caption= HX711\_morse\_experimantal Example Code,
label={lst:HX711MorseExample},
language=C++,
backgroundcolor=\color{backcolour},   
commentstyle=\color{codegreen},
keywordstyle=\color{magenta},
numberstyle=\tiny\color{codegray},
stringstyle=\color{codepurple},
basicstyle=\ttfamily\footnotesize,
breakatwhitespace=false,         
breaklines=true,                 
keepspaces=true,                 
numbers=left,       
numbersep=5pt,                  
showspaces=false,                
showstringspaces=false,
showtabs=false,                  
tabsize=2]{../../../Code/Syde/HTADemonstrator/Test/HTAOneLoadcell/HTAOneLoadcell.ino}

The sketch starts by including necessary libraries for Wi-Fi connectivity, web server management and communication with the HX711. The ESP32-S3 connects to a specified Wi-Fi network and sets up a web server on port 80. This server provides a webpage that guides users through three steps: taring the load cell, calibrating it with a known weight, and starting real-time weight measurements. The library functions \PYTHON{begin()}, \PYTHON{tare()}, \PYTHON{set\_scale()}, and \PYTHON{get\_units()} handle the initialization and measurement tasks. The \PYTHON{zero()} function uses \PYTHON{myScale.tare()} to reset the load cell to zero. The \PYTHON{enter()} function calculates a calibration factor by dividing the average weight reading by the known calibration weight, then applies it with \PYTHON{myScale.set\_scale()}. During operation, the \PYTHON{start() }function calls \PYTHON{myScale.get\_units(10)} to continuously measure and return the weight. The web interface, served from the \PYTHON{handleRoot() }function, provides three buttons: Zero, Enter, and Start. These buttons trigger API endpoints (\PYTHON{/zero}, \PYTHON{/enter}, and /start\PYTHON{/start}) that execute the corresponding functions on the ESP32-S3. The Start button also begins a loop that continuously updates the displayed weight value on the webpage by fetching data from the \PYTHON{/weight }endpoint every 100 milliseconds. Finally, serial data handling is implemented to allow debugging and status updates to be printed over USB. The program’s main \PYTHON{loop()} manages the web server and updates the web interface with real-time weight readings.


